Earnings calls, and especially the unscripted question-and-answer session, are a setting in which managers' language carries information beyond the scripted presentation \citep{matsumoto2011}. Disclosure theory holds that an informed manager may rationally withhold private information when disclosure is costly \citep{verrecchia1983}, and that non-disclosure can persist in equilibrium because outside investors cannot distinguish an uninformed manager from one who is informed but silent \citep{dye1985}, so a withholding state is itself informative. A firm privately committed to an acquisition it has not announced sits in a concrete such state: information about preliminary merger negotiations can be material well before any definitive agreement \citep{basic1988}, and although the firm may stay silent it may not mislead once it speaks \citep{rule10b5, basic1988}, so the executive can neither confirm nor deny the pending deal and is left to answer around it. When a manager cannot answer directly, the constraint can surface in the words themselves, for non-answers are common and not informationally empty: in a large share of calls managers leave requested information undisclosed, and, in the words of that literature, silence speaks \citep{hollander2010}.

Reading this language matters because the information in an undisclosed acquisition is exactly what securities law keeps off the call, yet the unscripted question-and-answer session is one of the few recurring venues in which a constrained executive must keep speaking and fielding questions. If the bind leaves a trace in the words, that trace would mark a deal's passage from private to public in the firm's own voice, in a channel separate from the price run-up already documented around mergers \citep{keown1981}.

The nearest prior work reads deal-related language along other dimensions. \citet{thewissen2024} show that acquirers manage the tone of their disclosure before stock-for-stock deals, and \citet{ragozzino2024} find that the volume of corporate-strategy vocabulary rises on the calls of acquisitive firms; both concern tone or subject matter rather than expressed uncertainty. Separately, \citet{keown1981} document that much of a merger's price reaction precedes its public announcement, an anticipatory signal that nonetheless lives in prices rather than in language. To our knowledge no prior work reads uncertainty language in the unscripted question-and-answer session in the anticipatory window before a withheld deal; we offer this as a positioning claim about where the contribution sits, not as a tested mechanism.

The purpose of this study is to ask whether chief-executive question-and-answer uncertainty tracks a deal's passage from private to public: whether it is elevated while an acquisition is withheld and recedes once it is announced. We construct a residual measure of chief-executive question-and-answer uncertainty, $\mathrm{UncResCEO}$, the call-varying component that remains once each executive's persistent speaking style is netted out, following a decomposition developed in the earnings-call language literature, and we defer its construction and the estimating equations to Section~2. We study how this measure behaves around cash acquisitions, using cash deals against a stock-acquisition comparison that shares the same withholding setting, on United States public, non-financial, non-utility firms' earnings calls from 2002 to 2018.

Our main findings can be previewed as follows. First, residual chief-executive uncertainty is elevated in the quarter before a cash acquisition, relative to the same firm's other quarters, with no comparable rise before stock acquisitions. Second, this elevation is anticipatory: it becomes indistinguishable from zero as soon as the deal is announced, even before it closes, whereas the cash that funds the purchase persists on the balance sheet until completion, two paths running on different clocks. This round-trip contrast is the study's strongest result.

Finally, the run-up concentrates in cash acquisitions rather than stock acquisitions, a difference that survives a formal pooled test of cash-specificity, though we read it as concentration rather than strict specificity, and the cash-accumulation mechanism behind it is left open. We also rule out the most immediate alternative. The run-up is not accounted for by analysts devoting more of the call to cash questions: with a direct measure of analyst cash scrutiny and its interaction entered alongside the pre-announcement indicator, the elevation survives and scrutiny does not amplify it. This means that scrutiny does not account for this particular run-up, not that scrutiny never matters. A further analysis turns to the firm's post-call information environment, measured by the bid-ask spread: the residual is unrelated to it, while the scripted presentation is positively associated with it, a pattern consistent with the scripted and unscripted segments reaching the outside environment in different ways.

This study makes four contributions, each descriptive. First, it reads a residual measure of chief-executive question-and-answer uncertainty, the call-varying component that remains once each executive's persistent speaking style is netted out, in the anticipatory window before an acquisition is public, where the nearest prior work has read managed tone, the volume of strategy vocabulary, and prices rather than uncertainty language. Second, it documents that this language pattern concentrates in cash acquisitions rather than stock, a contrast that survives a formal pooled test. Third, it reads the unscripted question-and-answer session as tracking a deal's passage from private to public, connecting the literature on when managers withhold private information with the evidence that markets anticipate corporate transactions, while claiming no causal channel. Fourth, it documents that this residual uncertainty is unrelated to the firm's post-call bid-ask spread while the scripted presentation is positively associated with it, consistent with a difference between the scripted and unscripted segments in how they relate to the outside information environment.

Throughout, we read these patterns as correlational and within-firm: the design supports no causal identification and establishes no mechanism. The concentration in cash deals is consistent with stock acquirers' incentive to manage their pre-deal narrative, an incentive that cash acquirers lack and that we develop in Section~2; we offer this as motivation, not identification. The source of the uncertainty, a compliance-constrained inability to speak versus a strategically chosen reticence, remains observationally equivalent.

The remainder of the paper proceeds as follows. Section~2 develops the conceptual framework and empirical strategy; Section~3 presents the three main analyses, the pre-announcement run-up, the differential timing around announcement, and the cash-specificity test; Section~4 reports additional analyses, ruling out analyst scrutiny, examining outsider reactions through the bid-ask spread, and a set of robustness checks---treating a withdrawal as a resolution event, dropping the cash equation's dynamic term, and re-running the main findings without the first-deal restriction; and Section~5 concludes.
